\section{Vòng lặp chung của một phép đo phản thực}
\label{sec:ch08-vong-lap-chung}

\index{đánh giá phản thực}%
\index{chỉ số đánh giá!ba lựa chọn tự do}%
Nếu không có bảng nào để chép thì phải đọc từng bài báo, và việc đọc ấy cho một
kết quả gọn hơn dự kiến. Các chỉ số dưới đây được đề xuất bởi những nhóm khác
nhau, ở thị giác máy tính và ở xử lý ngôn ngữ, dưới những cái tên không liên
quan gì nhau, nhưng chúng chạy chung một vòng lặp. Hình~\ref{fig:ch08-vong-lap-do} vẽ vòng lặp ấy.

\begin{figure}[htbp]
  \centering
  \begin{tikzpicture}[
  every node/.style={font=\footnotesize, align=center},
  box/.style={draw=black!65, rounded corners=2pt, fill=black!5,
    minimum width=22mm, minimum height=10mm},
  freebox/.style={box, dashed},
  flow/.style={-{Stealth[length=2mm]}, thick, black!70},
  note/.style={font=\scriptsize, align=center, black!55, text width=24mm}
]
  \node[box, minimum width=27mm] (model) {Mô hình $f$,\\đầu vào $x$};
  \node[box, right=5mm of model] (rank) {Xếp hạng\\đặc trưng};

  \node[freebox, below=14mm of rank] (subset) {Chọn tập con};
  \node[freebox, right=5mm of subset] (replace) {Thay giá trị};
  \node[box, right=5mm of replace] (rerun) {Chạy lại\\mô hình};
  \node[freebox, right=5mm of rerun] (compare) {So sánh\\đầu ra};

  \draw[flow] (model) -- (rank);
  \draw[flow] (rank) -- (subset);
  \draw[flow] (subset) -- (replace);
  \draw[flow] (replace) -- (rerun);
  \draw[flow] (rerun) -- (compare);

  \node[note, below=2mm of subset, text width=28mm]
    {$k$ điểm cao nhất, $k$ thấp nhất,\\hay ngẫu nhiên cỡ $n$};
  \node[note, below=2mm of replace] {số 0, trung bình,\\làm mờ, hay lấy mẫu};
  \node[note, below=2mm of compare, text width=34mm]
    {xác suất, độ chính xác,\\hay sau khi huấn luyện lại};
\end{tikzpicture}

  \caption{Vòng lặp mà mọi chỉ số trong chương này đều chạy. Ba ô viền đứt là
  ba chỗ mỗi chỉ số tự chọn lấy một cách làm; các họ chỉ số ở những mục sau khác
  nhau đúng ở ba ô đó và giống nhau ở phần còn lại. Đây là bản gộp của sách,
  dựng từ thủ tục mà từng bài báo mô tả, không phải sơ đồ của bài nào.}
  \label{fig:ch08-vong-lap-do}
\end{figure}

Vòng lặp đi như sau. Lời giải thích cho một thứ tự trên các đặc trưng của $x$.
Ta lấy một tập con theo thứ tự ấy, thay giá trị của các đặc trưng trong tập con
bằng một giá trị nào đó, chạy lại $f$ trên đầu vào đã sửa, rồi so đầu ra mới với
đầu ra cũ. Chênh lệch là con số của chỉ số. Lập luận đằng sau nó ngắn và có sức
thuyết phục: nếu lời giải thích đã chỉ đúng các đặc trưng mà $f$ dựa vào, thì bỏ
chúng đi phải làm $f$ đổi ý.

Ba ô viền đứt trong hình là ba chỗ vòng lặp không tự quyết định được, và cả
chương xoay quanh chúng.

\textbf{Chọn tập con nào.} Lấy $k$ đặc trưng điểm cao nhất là một cách; lấy $k$
đặc trưng điểm thấp nhất là một cách khác và đo một tính chất khác; lấy một tập
con ngẫu nhiên cỡ $n$ là cách thứ ba.

\textbf{Thay bằng giá trị gì.} Đặt về không, thay bằng trung bình của tập huấn
luyện, làm mờ, hay lấy mẫu từ một \tn{phân phối nền}{background distribution}. Chương~\ref{ch:shap} đã gặp câu
hỏi này khi SHAP phải định nghĩa thế nào là một đặc trưng vắng mặt, và
chương~\ref{ch:attribution-tu-nguyen-ly-dau} đã cho nó một tên tổng quát: chọn
giá trị thay thế chính là chọn \tn{độ đo}{measure} để lấy trung bình. Ở đây
câu hỏi cũ quay
lại ở một tầng mới. Nó không còn nằm trong phương pháp attribution nữa mà nằm
trong dụng cụ dùng để chấm điểm phương pháp ấy.

\textbf{So cái gì.} Xác suất của lớp được dự đoán, độ chính xác trên cả tập
kiểm tra, hay một đại lượng đo sau khi đã huấn luyện lại mô hình.

Ba lựa chọn ấy nhân với nhau, và mỗi tổ hợp là một chỉ số có tên riêng trong
tài liệu. Đó là lý do bốn mục tiếp theo xếp các chỉ số theo họ chứ không theo
tên: đọc theo tên thì chúng là một danh sách rời rạc, còn đọc theo ba ô viền đứt
thì chỉ có vài họ, và sự khác nhau giữa hai họ luôn quy về việc chúng điền gì
vào ba ô ấy.

Một hệ quả cần ghi ngay, vì nó là hình thức tổng quát của lời phê phán mà cả
Phần~IV sẽ triển khai. Không ô nào trong ba ô có một cách điền đúng theo nghĩa
toán học. Mỗi ô là một câu hỏi phản thực khác nhau, và
chương~\ref{ch:attribution-tu-nguyen-ly-dau} đã cho thấy hai câu hỏi phản thực
khác nhau có thể cùng hợp lệ mà cho hai kết quả khác nhau. Khi ấy hai chỉ số
xếp hạng ngược nhau không chứng tỏ một trong hai sai; chúng có thể đang đo hai
đại lượng khác nhau, cả hai đều được gọi là độ trung thực.
